\section{One more step: a single added lamp each round}\label{sec:one-more-step}

Section \ref{sec:wastlund} hands over the odd terms only. We change the doubling rule at one
point: after the lamps of a circle have been replaced by their images on the doubled circle,
add one lamp at the point of the doubled circle antipodal to $P$.

\begin{lemma}[the missing position]\label{lem:missing}
Let the $k$-th circle have circumference $C=2^{k+1}$, radius $2^{k}/\pi$, and lamps at the
arc positions $s=1,2,\dots,C-1$ measured from $P$. Under the replacement of Lemma
\ref{lem:doubling} the lamps of the doubled circle occupy every integer arc position
$1,\dots,2C-1$ except $s=C$, which is the point antipodal to $P$; and no lamp of the
$k$-th circle can be mapped there.
\end{lemma}

\begin{proof}
A lamp at arc position $s$ on the $k$-th circle is replaced by the two endpoints of the
diameter of the doubled circle through it, which stand at arc positions $s$ and $s+C$; this
covers each of $1,\dots,2C-1$ with $s\ne C$ exactly once. The point at arc position $C$ is the
far end of the diameter through $P$, that is, the line of centres; the other point of that
line on the $k$-th circle is $P$ itself, which carries no lamp. So the antipodal position has
no preimage.
\end{proof}

\begin{proposition}[brightness of the added lamps]\label{prop:added}
The lamp added at round $k$ stands at distance $2^{k+1}/\pi$ from $P$, hence contributes
$\pi^{2}/(2^{k+1})^{2}$ by \eqref{eq:inverse-square}, and the added lamps together give
\begin{equation}
  B=\frac{\pi^{2}}{2^{2}}+\frac{\pi^{2}}{4^{2}}+\frac{\pi^{2}}{8^{2}}+\cdots
   =\frac{\pi^{2}}{4}\cdot\frac{1}{1-\frac14}
   =\frac{\pi^{2}}{3}.
  \label{eq:supplement-brightness}
\end{equation}
Each term is the brightness at $P$ of one of the red lamps in Figure \ref{fig:onepic}, one term
per round, in the order in which the circles are drawn.
\end{proposition}

\begin{proof}
The distance from $P$ to the antipode is the diameter of the $k$-th circle, which is
$2^{k+1}/\pi$. The remaining sum is geometric with ratio $\frac14$.
\end{proof}

\begin{proof}[Proof of Theorem \ref{thm:direct}]
As the radius tends to infinity the circles flatten into the tangent line at $P$, chords tend
to arcs, and by Lemma \ref{lem:missing} every integer position is lit: the limiting line
carries lamps at all nonzero integers, at distances $1,2,3,\dots$ from $P$
(Figure \ref{fig:onepic}). The brightness they deliver at $P$ is
$\sum_{n\neq0}1/n^{2}=2S$. By Lemma \ref{lem:replacement} the replacement never changes the
brightness at $P$, so this total equals the brightness \eqref{eq:supplement-brightness} of the
lamps that had to be added. Hence
\begin{equation*}
  S=\frac{B}{2}
   =\frac{1}{2}\left(\frac{\pi^{2}}{2^{2}}+\frac{\pi^{2}}{4^{2}}+\frac{\pi^{2}}{8^{2}}+\cdots\right)
   =\frac{\pi^{2}}{6},
\end{equation*}
which is \eqref{eq:direct}.
\end{proof}

\begin{figure}[H]
\centering
\begin{tikzpicture}[scale=0.8, every node/.style={font=\small}]
  \coordinate (P) at (0,0);
  % Circles and line share one scale: one unit of arc length (or of distance along the
  % line) is 0.7 here, so the circle of circumference 2^{k+1} has radius 0.7*2^k/pi.
  \draw[thick] (0,0.223) circle (0.223);
  \draw[thick] (0,0.446) circle (0.446);
  \draw[thick] (0,0.891) circle (0.891);
  \draw[thick] (0,1.783) circle (1.783);
  \draw[thin,dashed] (0,0)--(0,3.7);
  % Lamps on each circle: the circle of circumference 2^{k+1} carries one lamp at every
  % arc distance s=1,...,2^{k+1}-1; the antipode s=2^k is the spot the red lamp fills in.
  % The innermost circle, of circumference 2, has only that antipode lamp.
  \foreach \s [evaluate=\s as \ang using {180*\s/2-90}] in {1,3}
    \fill[lampB] (0,0.446) + (\ang:0.446) circle (2.2pt);
  \foreach \s [evaluate=\s as \ang using {180*\s/4-90}] in {1,2,3,5,6,7}
    \fill[lampB] (0,0.891) + (\ang:0.891) circle (2.2pt);
  \foreach \s [evaluate=\s as \ang using {180*\s/8-90}] in {1,2,3,4,5,6,7,9,10,11,12,13,14,15}
    \fill[lampB] (0,1.783) + (\ang:1.783) circle (2.2pt);
  \fill[lampA] (0,0.446) circle (2.8pt);
  \fill[lampA] (0,0.891) circle (2.8pt);
  \fill[lampA] (0,1.783) circle (2.8pt);
  \fill[lampA] (0,3.565) circle (2.8pt);
  \draw[thick] (-5.1,0)--(5.1,0);
  \foreach \x/\lab in {-4.9/-7,-4.2/-6,-3.5/-5,-2.8/-4,-2.1/-3,-1.4/-2,-0.7/-1,0.7/1,1.4/2,2.1/3,2.8/4,3.5/5,4.2/6,4.9/7} {
    \fill[lampB] (\x,0) circle (2.2pt);
    \node[below,font=\footnotesize] at (\x,-0.05) {$\lab$};
  }
  \fill (P) circle (1.8pt);
  \node[above left] at (P) {$P$};
\end{tikzpicture}
\caption{Red marks the lamp added each round.}
\label{fig:onepic}
\end{figure}

The difference from Sections 2.4--2.5 is a single one. There every round is a pure equivalent
replacement, the total brightness is conserved, and the limit reads off the odd terms; here
the brightness accumulates round by round, the price is one added lamp per round, and what is
bought is that the series \eqref{eq:basel} itself appears on the figure. Each term of
\eqref{eq:supplement-brightness} can be pointed at on the picture. Once the picture is drawn
the only computation left is a geometric series with ratio $\frac14$.

\section*{Acknowledgments}

The geometric construction is due to Johan W\"astlund; this note only takes one further step
on his figure. The exposition also benefited from the video version by 3Blue1Brown.
